\documentclass[10pt,aspectratio=169]{beamer}

\usetheme{metropolis}
\usepackage{booktabs}
\usepackage{xcolor}
\usepackage{hyperref}
\usepackage{graphicx}
\usepackage{pgfpages}
\setbeameroption{show notes on second screen=right}

% ---------------------------------------------------------------------
% Yale colors (https://usebrand.yale.edu/colors)
% ---------------------------------------------------------------------
\definecolor{YaleBlue}{HTML}{00356B}
\definecolor{YaleBlueLight}{HTML}{286DC0}
\definecolor{YaleGray}{HTML}{4A4A4A}
\definecolor{YaleGrayLight}{HTML}{DDDDDD}
\definecolor{YaleAccent}{HTML}{BD5319}
\definecolor{goodgreen}{HTML}{2E8B57}

% ---------------------------------------------------------------------
% Metropolis theme overrides --- Yale Blue
% ---------------------------------------------------------------------
\setbeamercolor{normal text}{fg=black, bg=white}
\setbeamercolor{alerted text}{fg=YaleAccent}
\setbeamercolor{example text}{fg=YaleBlue}

\setbeamercolor{frametitle}{bg=YaleBlue, fg=white}
\setbeamercolor{title separator}{fg=YaleBlueLight}
\setbeamercolor{progress bar}{fg=YaleBlueLight, bg=YaleGrayLight}
\setbeamercolor{progress bar in section page}{fg=YaleBlueLight, bg=YaleGrayLight}

\setbeamercolor{section title}{fg=YaleBlue}
\setbeamercolor{subsection title}{fg=YaleBlueLight}
\setbeamercolor{standout}{bg=YaleBlue, fg=white}

\setbeamercolor{block title}{bg=YaleBlue, fg=white}
\setbeamercolor{block body}{bg=YaleGrayLight!30, fg=black}
\setbeamercolor{block title alerted}{bg=YaleAccent, fg=white}
\setbeamercolor{block body alerted}{bg=YaleGrayLight!30, fg=black}

\setbeamercolor{itemize item}{fg=YaleBlue}
\setbeamercolor{itemize subitem}{fg=YaleBlueLight}
\setbeamercolor{enumerate item}{fg=YaleBlue}

% ---------------------------------------------------------------------
% Meta
% ---------------------------------------------------------------------
\title{Macro News Pipeline}
\subtitle{Direct-link relevance framework}
% \author{James Zhang}
% \institute{Yale University}
\date{\today}

\begin{document}

\maketitle

\begin{frame}{Outline}
  \tableofcontents
\end{frame}
\section{The direct-link relevance rules}

% --- Rules overview ---
\begin{frame}{Rules for a valid asset mapping}
  \small
  An article--asset pair is relevant \textbf{only} when the asset is affected
  through one of five direct mechanisms. These rules are intentionally
  \emph{asset-class-agnostic} --- the same five cases apply to commodities,
  currencies, rates, equity indices, and sectors. In the future we may layer
  asset-class-specific rules on top (e.g., different thresholds for what
  counts as a ``dominant input'' for base metals vs.\ energy), but the
  current framework is general.

  \medskip
  \begin{enumerate}
    \item \textbf{NAMED} --- the asset appears in the text as a subject.
    \item \textbf{PRIMARY VEHICLE} --- the asset is the dominant traded
          instrument for a specific force described in the article.
    \item \textbf{ISSUING AUTHORITY} --- the article's subject is the
          institution that controls this asset's price.
    \item \textbf{SUBSTITUTION / INPUT LINK} --- a named commodity and this
          asset are closely-related substitutes or inputs.
    \item \textbf{SECTOR CONSTITUENT} --- a named company operates in this
          equity \emph{sector} (not index).
  \end{enumerate}

  \medskip
  The following slides walk through each rule with a concrete DJNW example.
\end{frame}

% --- Case 1 ---
\begin{frame}{Case 1: NAMED}
  \small
  \textbf{Rule.} The asset is named in the text as a \emph{subject} affected
  by a specific force --- not merely as a pricing unit, reporting currency,
  denominating proxy, or incidental mention.

  \medskip
  \begin{block}{Example: \textit{``ICCO Cocoa Daily Price --- Mar 21''}}
    The article reports the daily closing price of cocoa.
    Cocoa futures $\to$ \textbf{strong}, relevance 1.0. \\[4pt]
    The article originally also mentioned prices \textit{``in Euro/ton.''}
    Euro is a pricing unit here, not a subject --- rejected under case 1's
    exclusion clause. This led us to add an explicit
    \textit{functional-currency / reporting-currency} rejection rule after
    the same pattern leaked on the Carnival earnings article.
  \end{block}

  \medskip
  \textbf{Signal:} always \textbf{strong} (no reasoning chain needed).
\end{frame}

% --- Case 2 ---
\begin{frame}{Case 2: PRIMARY VEHICLE}
  \small
  \textbf{Rule.} The asset is the dominant traded instrument for the economic
  exposure described in the article. The exposure must be a \emph{specific
  force} --- not ``growth,'' ``uncertainty,'' or ``risk-off.'' The asset may
  be listed in a different jurisdiction than the named force.

  \medskip
  \begin{block}{Example: \textit{``China Jan-Feb Property Investment Rose 3.7\%''}}
    The article describes a slowdown in Chinese property construction.
    The \textbf{Hang Seng Index} is the primary listed market for taking
    direct equity exposure to the Chinese economy --- accepted as
    \textbf{weak 0.60}. \\[4pt]
    Without the cross-jurisdictional clause, the validator had rejected HSI
    by inventing ``the Hong Kong economy'' as a fake intermediate step.
    The safeguard we added stops the validator from fabricating intermediates.
  \end{block}

  \medskip
  \textbf{Signal:} always \textbf{weak} (the asset itself is not named).
\end{frame}

% --- Case 3 ---
\begin{frame}{Case 3: ISSUING AUTHORITY}
  \small
  \textbf{Rule.} The article's subject is the central bank, government
  agency, or supranational body that issues, controls, or directly sets
  this asset's price.

  \medskip
  \begin{block}{Example: \textit{``BOE's Rate Decision Could Benefit Sterling''}}
    The Bank of England is the issuing authority for UK monetary policy.
    \textbf{UK Long Gilt} $\to$ weak 0.60, \textbf{3-Month SONIA} $\to$
    weak 0.85. Sterling itself comes in via case 1 (named, strong 1.00).
    \\[4pt]
    Case 3 also drives the full UST curve on the
    \textit{``Dollar Firms, Ukraine War''} article, where the Fed's
    tightening stance is the named subject and the Fed is the issuing
    authority for US government debt.
  \end{block}

  \medskip
  \textbf{Signal:} always \textbf{weak} (the asset is controlled by the
  named institution, not named itself).
\end{frame}

% --- Case 4 ---
\begin{frame}{Case 4: SUBSTITUTION / INPUT LINK}
  \small
  \textbf{Rule.} The article names a commodity, input, or product, and this
  asset trades a closely-related substitute, input, or complement. The
  relationship need not be stated in the article --- well-known industry
  relationships count, as long as the original commodity is named.

  \medskip
  \begin{block}{Example: \textit{``China Jan-Feb Property Investment''} (cont.)}
    The article names \textit{``new construction starts.''} \textbf{LME
    Copper} (weak 0.50) and \textbf{LME Zinc} (weak 0.45) are accepted
    because Chinese construction is the dominant global demand driver for
    both metals. \\[4pt]
    LME Aluminium and LME Lead are \emph{rejected} on the same article ---
    property construction is not a dominant demand source for either.
    The guardrail is the word \emph{dominant}: ``any industrial use'' is
    not enough.
  \end{block}

  \medskip
  \textbf{Signal:} always \textbf{weak}.
  \textbf{Future work:} asset-class-specific rules may refine the
  ``dominant'' threshold for base metals vs.\ energy vs.\ agriculture.
\end{frame}

% --- Case 5 ---
\begin{frame}{Case 5: SECTOR CONSTITUENT}
  \small
  \textbf{Rule.} The article names a company or business event, and this
  asset is an equity \textbf{sector} (not an equity index) where the named
  company operates. Case 5 is restricted to sectors whose composition is
  defined by industry classification.

  \medskip
  \begin{block}{Example: \textit{``Carnival 1Q Loss/Shr \$1.66''}}
    Carnival Corporation is a major cruise-line operator.
    \textbf{Consumer Discretionary Sector} $\to$ weak 0.65. \\[4pt]
    Equity \emph{indices} (S\&P 500, Nasdaq 100, DAX, etc.) are
    excluded from case 5 because LLMs cannot reliably recall index
    composition. On the Russia/Boeing article, the model hallucinated
    that Boeing is a Nasdaq 100 constituent --- it is not.
    Indices can still be tagged via cases 1--4
    (e.g., HSI on China via case 2).
  \end{block}

  \medskip
  \textbf{Signal:} always \textbf{weak}. \\
  \textbf{Design decision:} restricting case 5 to sectors eliminates an
  entire class of factual-recall hallucinations at no recall cost.
\end{frame}

% --- Exclusion list ---
\begin{frame}{What does not count as a direct link}
  \scriptsize
  Each exclusion was motivated by a specific pattern the pipeline proposed
  that was judged incorrect upon review.

  \smallskip
  \begin{tabular}{@{}p{6.6cm}p{5.8cm}@{}}
    \toprule
    \textbf{Exclusion rule} & \textbf{Motivating example} \\
    \midrule
    Generic ``risk-off / safe-haven / global growth'' chains
      & Gold, JPY, Swiss franc tagged on every geopolitical article \\
    Historical correlations without article-specific pathway
      & ``Commodities tend to move with the dollar'' as standalone justification \\
    Secondary derivations through unnamed intermediates
      & China property $\to$ ``Australian economy'' $\to$ AUD \\
    Cross-country / cross-currency rate spillover
      & Ukraine/Dollar $\to$ Euro Bund, EURIBOR (Fed $\ne$ ECB) \\
    Functional-currency / reporting-currency boilerplate
      & Carnival 10-Q functional currency ``euro'' $\to$ Euro tagged \\
    Minority index/sector constituents
      & Small-cap company in a broad index proposed under case 5 \\
    Case 5 on equity indices (now sectors only)
      & Boeing $\to$ Nasdaq 100 (LLM hallucinated composition) \\
    Multi-step reasoning chains
      & Oil sanctions $\to$ ruble $\to$ emerging-market currencies \\
    Supply-chain-adjacent tiny exposure
      & China property $\to$ LME Aluminium \\
    \bottomrule
  \end{tabular}
\end{frame}

\begin{frame}{The exclusion list is a living document}
  \footnotesize
  Every rule on the previous slide started as a false positive the pipeline
  proposed and a human reviewer rejected. The process:

  \begin{enumerate}
    \item Run the pipeline on a new dataset.
    \item Inspect accepted and rejected tags with domain experts.
    \item When a pattern of incorrect proposals emerges, formalize it as an
          exclusion rule in the prompt.
    \item Re-run and verify the rule catches false positives without
          killing true positives.
  \end{enumerate}

  \medskip
  \textbf{Open question:} some exclusions may need asset-class-specific tuning:
  \begin{itemize}
    \item \textit{Substitution reach} (case 4) may need a stricter
          ``dominant demand driver'' threshold for base metals than energy.
    \item \textit{Cross-country rate spillover} may be too strict for
          tightly-coupled currency pairs (e.g., EUR/CHF).
  \end{itemize}

  \medskip
  For now, the rules are asset-class-agnostic.
\end{frame}

% ===================================================================
\section{Generator-verifier architecture}

\begin{frame}{Why split into mapper and validator}
  \small
  (Placeholder --- from ARCHITECTURE.md. Mapper recall-first, validator
  precision-first, base-rate shift between stages.)
\end{frame}

\begin{frame}{Cross-mapper agreement as a free reliability signal}
  \small
  (Placeholder --- the 70\% accept rate on \texttt{source=both} vs 13\% on
  single-source.)
\end{frame}

% ===================================================================
\section{Results on DJNW March 2022}

\begin{frame}{Aggregate numbers}
  \small
  (Placeholder --- accepted / rejected / strong / weak / source breakdown
  from results/direct\_djnw\_v2\_summary.json.)
\end{frame}

\begin{frame}{Focal articles walkthrough}
  \small
  (Placeholder --- BOE Sterling, USDA Soybeans, Russia aerospace, China property,
  Ukraine Dollar, Nymex Crude, ICCO Cocoa, ICE Sugar.)
\end{frame}

% ===================================================================
\section{Example mappings: before / after}

\begin{frame}{Reading these slides}
  \small
  Each example compares the mapping set produced at the last meeting against
  what the current pipeline produces on the same article, under the same
  direct-link rules.

  \medskip
  \begin{itemize}
    \item \textcolor{goodgreen}{\textbf{Green}}: tag that the current pipeline correctly adds or keeps.
    \item \textcolor{YaleAccent}{\textbf{Orange}}: tag correctly removed or still rejected.
    \item \textcolor{YaleGray}{\textbf{Gray}}: unchanged from last meeting.
  \end{itemize}

  \medskip
  The examples are chosen to show (a) the cases where iteration on the rules
  recovered real signal and (b) the cases where the stricter framework
  continues to reject chains it should reject.
\end{frame}

\begin{frame}{Example 1 --- China property investment (Jan-Feb)}
  \footnotesize
  \textit{``China Jan-Feb Property Investment Rose 3.7\% on Year''} --- slowdown in
  new construction starts; Hong Kong property developers mentioned.

  \medskip
  \begin{tabular}{@{}lll@{}}
    \toprule
    \textbf{Asset} & \textbf{Last meeting} & \textbf{Now} \\
    \midrule
    Hang Seng Index        & --          & \textcolor{goodgreen}{weak 0.60 \checkmark} \\
    LME Copper             & --          & \textcolor{goodgreen}{weak 0.50 \checkmark} \\
    Copper (High Grade)    & --          & \textcolor{goodgreen}{weak 0.40 \checkmark} \\
    LME Zinc               & --          & \textcolor{goodgreen}{weak 0.45 \checkmark} \\
    LME Aluminium          & --          & \textcolor{YaleAccent}{rejected} \\
    LME Lead               & --          & \textcolor{YaleAccent}{rejected} \\
    Australian Dollar      & --          & \textcolor{YaleAccent}{rejected} \\
    \bottomrule
  \end{tabular}

  \medskip
  HSI returns via case 2 (primary listed market for China equity exposure).
  Copper and Zinc come in via case 4 (China construction is the dominant global
  demand driver for both metals).
\end{frame}

\begin{frame}{Why copper and zinc are a real case 4, not a floodgate}
  \small
  \textbf{Recommendation:} tag LME Copper / Copper HG / LME Zinc on this article.

  \medskip
  \begin{itemize}
    \item The article explicitly names \textit{``new construction starts''} and
          \textit{``real estate investment in China''} --- not a generic growth reference.
    \item Chinese construction is the \emph{single largest} demand source for
          both copper (wire, plumbing) and zinc (galvanized steel). This is the
          textbook case-4 substitution/input link.
    \item The guardrail is holding: on the same article, LME Aluminium, LME
          Lead, and Australian Dollar are rejected as too indirect.
    \item Aggregate commodity FPR \emph{fell} between versions (from 53\% to
          43\%). If case 4 were over-firing, this would move the other way.
  \end{itemize}

  \medskip
  The rule we should hold the line on: \emph{``dominant demand driver,'' not
  ``any industrial use.''}
\end{frame}

\begin{frame}{Example 2 --- Dollar firms on Ukraine war}
  \footnotesize
  \textit{``Dollar Firms, Ukraine War Should Send It Higher''} --- USD strength,
  Fed in tightening mode, safe-haven flows into US assets.

  \medskip
  \begin{tabular}{@{}lll@{}}
    \toprule
    \textbf{Asset} & \textbf{Last meeting} & \textbf{Now} \\
    \midrule
    Euro                         & \textcolor{YaleGray}{strong 0.85} & strong 0.85 \\
    Japanese Yen                 & \textcolor{YaleGray}{strong 0.75} & strong 0.75 \\
    3-Month SOFR                 & \textcolor{YaleGray}{weak 0.70}   & weak 0.70 \\
    10 Year US Government Note   & \textcolor{YaleGray}{weak 0.75}   & weak 0.75 \\
    2/5/30Y UST, Ultra T-Bond    & --                                & \textcolor{goodgreen}{weak 0.50--0.65 \checkmark} \\
    \midrule
    Euro Bund / Bobl             & --  & \textcolor{YaleAccent}{rejected} \\
    3-Month EURIBOR              & --  & \textcolor{YaleAccent}{rejected} \\
    DAX Index                    & --  & \textcolor{YaleAccent}{rejected} \\
    Swedish Krona                & --  & \textcolor{YaleAccent}{rejected} \\
    \bottomrule
  \end{tabular}

  \medskip
  The full UST curve comes in via case 3 (Fed is the issuing authority for
  USTs, and the article explicitly names the Fed's policy stance). Cross-country
  rate spillover rejections still hold cleanly.
\end{frame}

\begin{frame}{Example 3 --- BOE rate decision}
  \footnotesize
  \textit{``BOE's Rate Decision Could Benefit Sterling in Short-Term''}

  \medskip
  \begin{tabular}{@{}lll@{}}
    \toprule
    \textbf{Asset} & \textbf{Last meeting} & \textbf{Now} \\
    \midrule
    Sterling        & \textcolor{YaleGray}{strong 1.00} & strong 1.00 \\
    UK Long Gilt    & \textcolor{YaleGray}{weak 0.60}   & weak 0.60 \\
    3-Month SONIA   & \textcolor{YaleGray}{weak 0.85}   & weak 0.85 \\
    \bottomrule
  \end{tabular}

  \medskip
  Clean case 1 (Sterling named) plus case 3 (BOE is the issuing authority for
  gilts and SONIA). No change across versions --- this is what a stable focal
  article looks like.
\end{frame}

\begin{frame}{Example 4 --- The Carnival functional-currency leak}
  \footnotesize
  \textit{``Carnival 1Q Loss/Shr \$1.66''} --- earnings note; the 10-Q risk-factor
  section lists Euro and Sterling as functional currencies.

  \medskip
  \begin{tabular}{@{}lll@{}}
    \toprule
    \textbf{Asset} & \textbf{Last meeting} & \textbf{Now (pre-fix)} & \textbf{Now (fixed)} \\
    \midrule
    Euro      & strong 0.20 & strong 0.20 & \textcolor{YaleAccent}{rejected} \\
    Sterling  & weak 0.20   & weak 0.20   & \textcolor{YaleAccent}{rejected} \\
    Consumer Discretionary & \textcolor{YaleGray}{weak 0.65} & weak 0.65 & weak 0.65 \\
    Brent / WTI / Gas Oil & \textcolor{YaleGray}{weak 0.45} & weak 0.45 & weak 0.45 \\
    \bottomrule
  \end{tabular}

  \medskip
  The validator was literal-matching ``euro'' in the risk-factor section
  (``subject to foreign currency translational risk'') and upgrading to strong.
  Case 1 already excluded ``reporting currency / denominating proxy,'' but not
  prominently enough; added an explicit rejection rule to both mapper and
  validator prompts.
\end{frame}

\begin{frame}{What the examples collectively show}
  \small
  \begin{itemize}
    \item \textbf{Cross-jurisdictional case 2} recovered HSI on China without
          breaking anything else --- validator chain-rejection safeguard prevents
          it from inventing ``host country economy'' intermediates.
    \item \textbf{Case 4} tags are well-anchored: copper/zinc on China
          construction pass; aluminium, lead, AUD do not. Guardrail is the
          word \emph{dominant}.
    \item \textbf{Case 3} (issuing authority) is doing real work on rate stories
          --- full UST curve on Fed discussion, Gilt + SONIA on BOE.
    \item \textbf{Stable focals} (BOE, Cocoa, Sugar, Nymex Crude, Soybeans)
          are unchanged, which is the expected behavior for well-specified
          case-1 stories.
    \item \textbf{Remaining leak}: functional-currency boilerplate upgraded to
          case 1. Fix is a one-line prompt addition to both mapper and
          validator.
  \end{itemize}
\end{frame}

% ===================================================================
\section{Open design questions}

\begin{frame}{Where the framework has judgment calls}
  \small
  (Placeholder --- e.g. asset universe trim, base metals on China property,
  substitution reach, validator-only vs two-mapper architecture, next dataset.)
\end{frame}

\end{document}
